\chapter{Densidad y numerabilidad}\label{chap:06}
\section{Aproximar un espacio con pocos puntos}
Una sucesión solo utiliza un conjunto numerable de términos. Es natural preguntar cuándo todo un espacio puede describirse, en términos de aproximación, a partir de un conjunto numerable. Este capítulo desarrolla la lectura 6 de \cite{valencia}: densidad, separabilidad, bases numerables y el principio de Lindelöf. Se utilizarán las propiedades de clausura, subespacios y continuidad ya establecidas.

Un conjunto es \emph{a lo sumo numerable} si es finito o puede enumerarse mediante números naturales. Usaremos que una unión numerable de conjuntos numerables es numerable y que un producto finito de conjuntos numerables es numerable. Por ejemplo, los pares de naturales pueden enumerarse recorriendo las diagonales $m+n=2,3,\ldots$; una familia numerable de enumeraciones se reduce así a una lista de pares. Las elecciones necesarias se entienden en el marco usual de la teoría de conjuntos con el axioma de elección.

\begin{teorema}\label{ch06:denso}\index{conjunto!denso}
Para $A\subset X$, con $X$ métrico no vacío, son equivalentes:
\begin{enumerate}[label=(\roman*)]
\item $A\ne\varnothing$ y $d(x,A)=0$ para todo $x\in X$.
\item Todo abierto no vacío de $X$ intersecta a $A$.
\item $\cl A=X$.
\end{enumerate}
En este caso se dice que $A$ es denso en $X$.
\end{teorema}
\begin{proof}
Si (i) se cumple y $U$ es abierto no vacío, elijamos $x\in U$ y $r>0$ con $B(x,r)\subset U$. La igualdad $d(x,A)=0$ produce $a\in A$ con $d(x,a)<r$, así que $a\in U$. Si (ii) se cumple, cada bola $B(x,r)$ es un abierto no vacío e intersecta a $A$, lo que implica $x\in\cl A$ para cada $x$. Finalmente, si $\cl A=X$, entonces $A$ es no vacío porque $X$ lo es, y el ejercicio 7 del capítulo~\ref{chap:03} da $d(x,A)=0$.
\end{proof}

\begin{ejemplo}[Racionales y truncaciones]
En un espacio discreto, cada singleton es abierto. Por tanto un conjunto denso debe contener cada punto, y el único conjunto denso es el espacio entero.

En $\R$, $\Q$ es denso. Si $a<b$, elegimos $n$ con $1/n<b-a$. Existe un entero $m$ tal que $na<m\le na+1$; por ejemplo, $m=\lfloor na\rfloor+1$. Entonces $a<m/n\le a+1/n<b$. Así cada intervalo abierto no vacío contiene un racional.

En $\ell^1$, el subespacio $c_{00}=\operatorname{span}\{e_k:k\ge1\}$ de sucesiones con soporte finito es denso. Para $a\in\ell^1$, su truncación $a^{[N]}=(a_1,\ldots,a_N,0,\ldots)$ satisface
\[
\norm{a-a^{[N]}}_1=\sum_{k>N}|a_k|\longrightarrow0.
\]
La convergencia de la serie justifica que sus colas tiendan a cero. Sin embargo, $c_{00}\ne\ell^1$: la sucesión $(1/n^2)$ pertenece a $\ell^1$ y tiene infinitas coordenadas no nulas. Una familia puede generar un subespacio \emph{denso} sin generar algebraicamente todo el espacio.\index{c00@$c_{00}$}
\end{ejemplo}
Más adelante se demostrará que los polinomios son densos en $C([0,1])$ con la norma uniforme. Esa afirmación necesita un teorema de aproximación, y no se usará como premisa antes de demostrarlo.

\begin{definicion}\index{separabilidad}
Un espacio métrico es separable si posee un subconjunto denso a lo sumo numerable.
\end{definicion}
Así, $\R$ es separable aunque no sea numerable. Un espacio discreto es separable exactamente cuando es numerable, porque su único subconjunto denso es él mismo. No debe confundirse la cantidad de puntos del espacio con la cantidad de puntos necesaria para aproximarlos.

\section{Bases numerables}
\begin{definicion}\index{base!de una topologia@de una topología}\index{segunda numerabilidad}
Una base de la topología de $X$ es una familia $\mathcal B$ de abiertos tal que todo abierto de $X$ es unión de elementos de $\mathcal B$. El espacio es segundo numerable si posee una base a lo sumo numerable.
\end{definicion}
Una formulación equivalente, útil para las pruebas, es: para cada abierto $U$ y cada $x\in U$, existe $B\in\mathcal B$ con $x\in B\subset U$. Si todo abierto es unión de elementos de la base, uno de los que componen $U$ contiene a $x$. En sentido contrario, tomando esos elementos para todos los $x\in U$, su unión es exactamente $U$.

\begin{lema}\label{ch06:baseadenso}
De una base $\mathcal B$ puede obtenerse un conjunto denso cuya cardinalidad no exceda la de $\mathcal B$.
\end{lema}
\begin{proof}
Eliminemos el conjunto vacío de la base, si aparece: no aporta puntos a ninguna unión, por lo que esta sigue siendo base. Para cada $B\in\mathcal B$ elegimos $x_B\in B$, y ponemos $D=\{x_B:B\in\mathcal B\}$. Todo abierto no vacío contiene algún elemento no vacío de la base y, en consecuencia, contiene el punto elegido en él. Por el teorema~\ref{ch06:denso}, $D$ es denso. La aplicación $B\mapsto x_B$ es una sobreyección de la base sobre $D$, lo que da la cota cardinal.
\end{proof}

\begin{teorema}\label{ch06:separablebase}
Un espacio métrico es separable si y solo si es segundo numerable.
\end{teorema}
\begin{proof}
Una base numerable proporciona un conjunto denso numerable por el lema anterior. Para la otra dirección, sea $D=\{a_n:n\in\N\}$ denso; si $D$ es finito no vacío, se pueden repetir sus elementos en una enumeración. Definamos
\[
\mathcal B=\{B(a_n,2^{-k}):n,k\in\N\}.
\]
Esta familia es numerable porque está indexada por pares de naturales. Sea $U$ abierto y $x\in U$; elijamos $\eps>0$ con $B(x,\eps)\subset U$, y luego $k$ con $2^{1-k}<\eps$. La densidad permite elegir $a_n$ con $d(x,a_n)<2^{-k}$. Entonces $x\in B(a_n,2^{-k})$ y, para $y$ en esa bola,
\[
d(x,y)\le d(x,a_n)+d(a_n,y)<2\cdot2^{-k}<\eps.
\]
Por tanto $x\in B(a_n,2^{-k})\subset U$. Esta es la propiedad de base. En el espacio vacío, la familia vacía es una base y el conjunto vacío es denso en él, así que ambas propiedades también coinciden.
\end{proof}
La prueba revela por qué los centros densos son suficientes: se aproxima primero el centro deseado, reservando la mitad de la tolerancia para el radio de la bola. No basta elegir un centro cercano sin controlar también ese radio.

\begin{teorema}[Lindelöf]\label{ch06:lindelof}\index{Lindelof@Lindelöf, teorema de}
Sea $X$ métrico separable y $\mathcal U$ una familia de abiertos. Existe una subfamilia a lo sumo numerable con la misma unión que $\mathcal U$. En particular, toda cubierta abierta de $X$ admite una subcubierta numerable.
\end{teorema}
\begin{proof}
Sea $\mathcal B=\{B_n:n\in\N\}$ una base numerable. Consideremos los índices $n$ para los que $B_n$ está contenido en algún miembro de $\mathcal U$, y para cada uno elijamos un solo $U_n\in\mathcal U$ con $B_n\subset U_n$. La familia elegida es numerable.

Si $x\in\bigcup\mathcal U$, existe $U\in\mathcal U$ con $x\in U$. La propiedad de base proporciona un $B_n$ con $x\in B_n\subset U$. Este índice está entre los seleccionados; aunque el abierto elegido $U_n$ pueda ser distinto de $U$, contiene a $B_n$ y, por tanto, a $x$. Así $\bigcup\mathcal U\subset\bigcup_nU_n$. La inclusión contraria se sigue de que todos los $U_n$ pertenecen a $\mathcal U$. Si la unión es vacía, la subfamilia vacía sirve.
\end{proof}
El teorema no dice que la subcubierta pueda ser finita; esa exigencia define una propiedad más fuerte que estudiaremos bajo el nombre de compacidad.

\section{Ejercicios resueltos}
\begin{ejercicio}[Lectura 6, ejercicio 1]
Demostrar que $\ell^\infty$ no es separable.
\end{ejercicio}
\begin{solucion}
Sea $E=\{0,1\}^{\N}\subset\ell^\infty$. Este conjunto no es numerable: si sus elementos se listaran como $a^{(1)},a^{(2)},\ldots$, la sucesión $b$ definida por $b_n=1-a^{(n)}_n$ diferiría de la sucesión número $n$ en la coordenada $n$ y no aparecería en la lista. Es el argumento diagonal de Cantor.

Para $a,b\in E$ distintos, existe una coordenada donde difieren en $1$, y ninguna diferencia tiene módulo mayor que $1$; por tanto $\norm{a-b}_\infty=1$. Las bolas $B(a,1/3)$, $a\in E$, son disjuntas dos a dos: un punto común daría por la triangular una distancia entre centros menor que $2/3$.

Si $D=\{d_n:n\in\N\}$ fuera denso, cada bola contendría algún $d_n$. Asignemos a cada $a\in E$ el primer índice de un punto de $D$ que pertenece a $B(a,1/3)$. Bolas disjuntas reciben índices distintos, lo que produciría una inyección del conjunto no numerable $E$ en $\N$. Es imposible; por tanto $\ell^\infty$ no es separable.
\end{solucion}

\begin{ejercicio}[Lectura 6, ejercicio 2]
Demostrar que un espacio normado $V$ es separable si y solo si existe un conjunto a lo sumo numerable $L\subset V$ cuyo subespacio generado sea denso.
\end{ejercicio}
\begin{solucion}
Si $D$ es denso y numerable, podemos tomar $L=D$, pues $D\subset\operatorname{span}D$ implica densidad del subespacio generado.

Para el recíproco, sea $\F_0=\Q$ en el caso real y $\F_0=\Q+i\Q$ en el complejo. Sea $D$ el conjunto de todas las combinaciones lineales finitas de elementos de $L$ con coeficientes en $\F_0$, incluyendo el vector cero. Para cada longitud de combinación hay solo un producto finito de conjuntos numerables; al tomar la unión sobre todas las longitudes, $D$ es numerable.

Probemos su densidad. Dado $v\in V$ y $\eps>0$, la densidad de $\operatorname{span}L$ permite escoger $w=\sum_{j=1}^m\alpha_jv_j$ con $\norm{v-w}<\eps/2$. Para cada $j$, elijamos $\beta_j\in\F_0$ tal que
\[
|\alpha_j-\beta_j|<\frac{\eps}{2m(1+\norm{v_j})}.
\]
Los racionales son densos en $\R$, y $\Q+i\Q$ es denso en $\C$ al aproximar por separado las partes real e imaginaria; por ello estas elecciones son posibles. Entonces $u=\sum_{j=1}^m\beta_jv_j\in D$ y
\[
\norm{w-u}\le\sum_{j=1}^m|\alpha_j-\beta_j|\norm{v_j}<\eps/2.
\]
La triangular da $\norm{v-u}<\eps$. Si la combinación es vacía, se toma $u=0$; en particular, $V=\{0\}$ queda cubierto. Así $D$ es denso.
\end{solucion}

\begin{ejercicio}[Lectura 6, ejercicio 3]
Si $(A_n)$ es una familia a lo sumo numerable de subespacios separables de un espacio métrico $X$, demostrar que $Y=\bigcup_nA_n$ es separable.
\end{ejercicio}
\begin{solucion}
Para cada $A_n$ no vacío, elijamos $D_n\subset A_n$ denso en su topología relativa y numerable; para el vacío tomemos $D_n=\varnothing$. El conjunto $D=\bigcup_nD_n$ es numerable. Si $y\in Y$ y $r>0$, existe $n$ con $y\in A_n$. Por densidad relativa, $B_X(y,r)\cap A_n$ contiene un punto de $D_n$, y por tanto $B_Y(y,r)$ contiene un punto de $D$. Esto demuestra $\cl D^{\,Y}=Y$. No se necesita que los $A_n$ sean abiertos ni cerrados.
\end{solucion}

\begin{ejercicio}[Lectura 6, ejercicio 4]
Demostrar que todo subespacio de un espacio métrico separable es separable.
\end{ejercicio}
\begin{solucion}
Sea $\{B_n\}$ una base numerable de $X$, cuya existencia asegura el teorema~\ref{ch06:separablebase}. Los conjuntos $Y\cap B_n$ forman una base de $Y$: si $U=Y\cap G$ es abierto relativo y $y\in U$, hay $B_n$ con $y\in B_n\subset G$, y entonces $y\in Y\cap B_n\subset U$. Esta base es numerable. Elegimos un punto de cada elemento no vacío de ella; por el lema~\ref{ch06:baseadenso}, esos puntos forman un denso numerable de $Y$.

No basta intersectar con $Y$ un denso numerable prefijado de $X$. Por ejemplo, $\Q$ es denso en $\R$, pero su intersección con $Y=\R\setminus\Q$ es vacía. La prueba con una base elige puntos dentro del subespacio y evita ese error.
\end{solucion}

\begin{ejercicio}[Lectura 6, ejercicio 5]
Sean $X_1,\ldots,X_m$ espacios métricos no vacíos. Demostrar que su producto finito, con una métrica producto, es separable si y solo si lo es cada factor.
\end{ejercicio}
\begin{solucion}
Podemos usar $d(x,y)=\max_{1\le j\le m}d_j(x_j,y_j)$. Su bola de radio $r$ es el producto de las bolas de radio $r$ de los factores. Si $D_j$ es denso numerable en $X_j$, el producto $D=\prod_{j=1}^mD_j$ es numerable. Dado $x$ y $r>0$, elegimos $a_j\in D_j$ con $d_j(x_j,a_j)<r$ para cada uno de los finitísimos índices. Entonces $a\in D\cap B(x,r)$ y $D$ es denso.

Recíprocamente, sea $D$ denso numerable en el producto. La proyección $\pi_j(D)$ es numerable. Dado $x_j\in X_j$ y $r>0$, elegimos puntos arbitrarios en los demás factores y formamos $x$ con coordenada $j$ igual a $x_j$. Una aproximación $a\in D$ con $d(a,x)<r$ proporciona $d_j(a_j,x_j)<r$, así que $\pi_j(D)$ es denso en $X_j$.

Las métricas producto finitas dadas por suma, máximo o norma $p$ inducen la misma topología: para números no negativos $t_j$, $\max_jt_j\le(\sum_jt_j^p)^{1/p}\le m^{1/p}\max_jt_j$ y $\max_jt_j\le\sum_jt_j\le m\max_jt_j$. Estas cotas comparan las bolas y trasladan el resultado a cualquiera de esas métricas. La hipótesis de factores no vacíos es necesaria: un factor vacío hace que el producto sea vacío, aunque otro factor sea no separable.
\end{solucion}

\begin{definicion}\index{denso en ninguna parte}
Un conjunto $A\subset X$ es denso en ninguna parte si $\operatorname{int}(\cl A)=\varnothing$. La clausura dentro de esta definición es esencial: tener interior vacío por sí solo no basta, como muestra $\Q\subset\R$.
\end{definicion}

\begin{ejercicio}[Lectura 6, ejercicio 6]
Demostrar que un conjunto denso en ninguna parte tiene complemento denso, y que un conjunto denso y cerrado tiene complemento denso en ninguna parte.
\end{ejercicio}
\begin{solucion}
Si $\operatorname{int}(\cl A)=\varnothing$, ningún abierto no vacío $U$ está contenido en $\cl A$. Por tanto $U\cap(X\setminus\cl A)\ne\varnothing$, y en particular $U\cap(X\setminus A)\ne\varnothing$. Así el complemento de $A$ es denso. Hemos probado algo más fuerte: el complemento de la clausura es abierto y denso.

Si $A$ es denso y cerrado, $A=\cl A=X$; por tanto su complemento es vacío, cuya clausura tiene interior vacío. La segunda afirmación de la lectura es correcta, aunque trivial con ambas hipótesis. Una formulación adicional útil es que el complemento de un abierto denso es cerrado y denso en ninguna parte: si $G$ es abierto y denso, $F=X\setminus G$ es cerrado y no contiene ningún abierto no vacío, porque todo abierto no vacío intersecta a $G$. Así $\operatorname{int}(\cl F)=\operatorname{int}F=\varnothing$.
\end{solucion}

\begin{ejercicio}[Lectura 6, ejercicio 7]
Demostrar que en un espacio métrico no vacío un conjunto no puede ser simultáneamente denso y denso en ninguna parte.
\end{ejercicio}
\begin{solucion}
Si $A$ es denso, $\cl A=X$. Entonces $\operatorname{int}(\cl A)=\operatorname{int}X=X\ne\varnothing$, lo que contradice ser denso en ninguna parte. La condición de espacio no vacío evita la excepción $X=A=\varnothing$.
\end{solucion}

\begin{ejercicio}[Lectura 6, ejercicio 8]\label{ch06:intervalos}
Demostrar que todo abierto de $\R$ es unión a lo sumo numerable de intervalos abiertos disjuntos dos a dos.
\end{ejercicio}
\begin{solucion}
Sea $U\subset\R$ abierto. Si es vacío, la unión vacía resuelve el problema. Para $x\in U$, definamos $I_x$ como la unión de todos los intervalos abiertos contenidos en $U$ que contienen a $x$. Hay al menos uno porque $U$ es abierto. La unión es abierta y está contenida en $U$.

Además, $I_x$ es un intervalo: si $a<b$ pertenecen a $I_x$, hay intervalos $J_a,J_b$ de la familia que contienen respectivamente a $a,b$ y ambos contienen a $x$. Como se intersectan, su unión es un intervalo. Para justificar este último hecho, un punto entre dos puntos de $J_a\cup J_b$ queda en uno de los intervalos, o entre ellos; en el segundo caso un punto común a ambos permite dividir el segmento en dos partes contenidas en los respectivos intervalos. Por ello todos los puntos entre $a,b$ pertenecen a $J_a\cup J_b\subset I_x$.

Así $I_x$ es un intervalo abierto, posiblemente una semirrecta o toda $\R$. Es maximal entre los intervalos de $U$ que contienen a $x$, porque su definición ya reúne a todos ellos. Si $I_x$ e $I_y$ se intersectan, su unión es un intervalo abierto contenido en $U$; contiene tanto a $x$ como a $y$. Por maximalidad está contenida tanto en $I_x$ como en $I_y$, y por tanto $I_x=I_y$. Los intervalos distintos de esta colección son disjuntos y cubren $U$ porque $x\in I_x$.

Cada intervalo no vacío contiene un racional. Fijemos una enumeración de $\Q$ y asignemos a cada intervalo el primer racional de la lista que contiene. Dos intervalos disjuntos reciben racionales distintos, de modo que la familia se inyecta en $\Q$ y es a lo sumo numerable. Esto concluye la descomposición, incluidos los intervalos no acotados.
\end{solucion}

\section{La utilidad de una base numerable}
En un espacio métrico, una familia numerable de puntos de aproximación se convierte en una familia numerable de entornos básicos. Esa conversión permite reducir cubiertas abiertas a subfamilias numerables y heredar separabilidad a todos los subespacios. La condición «denso en ninguna parte» mide una limitación distinta: ni siquiera después de agregar puntos límite aparece una región abierta completa.
